\section{Homework Week 8}

\begin{exercise}
    Let \(A\) be a finite-dimensional algebra over a field \(k\).
    Suppose that \(V\) is a finitely generated \(A\)-module, and that a certain simple
    \(A\)-module \(S\) occurs as a composition factor of \(V\) with multiplicity \(1\).
    Suppose that there exist nonzero homomorphisms
    \[
    S \longrightarrow V
    \qquad\text{and}\qquad
    V \longrightarrow S.
    \]
    Prove that \(S\) is a direct summand of \(V\).
    \end{exercise}
    
    \begin{proof}
    Since \(A\) is finite-dimensional over \(k\) and \(V\) is finitely generated over \(A\),
    the module \(V\) is finite-dimensional over \(k\). Hence \(V\) has finite length.
    
    Let
    \[
    \iota:S\longrightarrow V,
    \qquad
    \rho:V\longrightarrow S
    \]
    be nonzero \(A\)-homomorphisms. Since \(S\) is simple, \(\iota\) is injective and
    \(\rho\) is surjective.
    
    We claim that
    \[
    \rho\circ \iota\neq 0.
    \]
    Suppose otherwise that \(\rho\circ\iota=0\). Then
    \[
    \iota(S)\subseteq \ker\rho.
    \]
    Thus \(S\) occurs as a composition factor of \(\ker\rho\). Since \(\rho\) is surjective,
    we have a short exact sequence
    \[
    0\longrightarrow \ker\rho \longrightarrow V \longrightarrow S \longrightarrow 0.
    \]
    By additivity of composition multiplicities in short exact sequences,
    \[
    [V:S]=[\ker\rho:S]+[S:S]\geq 1+1=2,
    \]
    contradicting the assumption that \(S\) occurs in \(V\) with multiplicity \(1\).
    Therefore
    \[
    \rho\circ\iota\neq 0.
    \]
    
    By Schur's lemma, the nonzero endomorphism
    \[
    \rho\circ\iota:S\longrightarrow S
    \]
    is an isomorphism. Define
    \[
    \sigma=\iota\circ(\rho\circ\iota)^{-1}:S\longrightarrow V.
    \]
    Then
    \[
    \rho\circ\sigma
    =
    \rho\circ\iota\circ(\rho\circ\iota)^{-1}
    =
    \operatorname{id}_S.
    \]
    Hence \(\rho:V\to S\) splits. Therefore
    \[
    V=\ker\rho\oplus \sigma(S),
    \]
    and since \(\sigma(S)\cong S\), the simple module \(S\) is a direct summand of \(V\).
    \end{proof}
    
    
    \begin{exercise}
    Suppose that \(U,V,W\) are modules over a finite-dimensional algebra \(A\) over a field,
    and let
    \[
    p_W:P_W\longrightarrow W
    \]
    be a projective cover of \(W\). Assume that all modules involved are finite-dimensional.
    
    \begin{enumerate}
        \item[(a)] Show that a homomorphism
        \[
        \alpha:U\longrightarrow W
        \]
        is an essential epimorphism if and only if there is a surjective homomorphism
        \[
        h:P_W\longrightarrow U
        \]
        such that the composite
        \[
        P_W \xrightarrow{\,h\,} U \xrightarrow{\,\alpha\,} W
        \]
        is a projective cover of \(W\). In this situation, show that
        \[
        h:P_W\longrightarrow U
        \]
        must be a projective cover of \(U\).
    
        \item[(b)] Prove the following ``extension and converse'' to Nakayama's lemma:
        let \(V\) be any submodule of \(U\). Then the canonical epimorphism
        \[
        q:U\longrightarrow U/V
        \]
        is an essential epimorphism if and only if
        \[
        V\subseteq \operatorname{Rad}U.
        \]
    \end{enumerate}
    \end{exercise}
    
    \begin{proof}
    \textbf{(a)}
    Recall that an epimorphism \(f:X\to Y\) is essential if whenever \(X_0\leq X\) satisfies
    \(f(X_0)=Y\), one must have \(X_0=X\).
    
    First suppose that
    \[
    \alpha:U\longrightarrow W
    \]
    is an essential epimorphism. Since \(P_W\) is projective and \(\alpha\) is surjective,
    the map \(p_W:P_W\to W\) lifts through \(\alpha\). Thus there exists an \(A\)-homomorphism
    \[
    h:P_W\longrightarrow U
    \]
    such that
    \[
    \alpha\circ h=p_W.
    \]
    Because \(p_W\) is surjective, we have
    \[
    \alpha(h(P_W))=W.
    \]
    Since \(\alpha\) is essential, this forces
    \[
    h(P_W)=U.
    \]
    Hence \(h\) is surjective. Moreover, the composite
    \[
    P_W \xrightarrow{\,h\,} U \xrightarrow{\,\alpha\,} W
    \]
    is precisely \(p_W\), hence is a projective cover of \(W\).
    
    Conversely, suppose that there exists a surjective homomorphism
    \[
    h:P_W\longrightarrow U
    \]
    such that
    \[
    \alpha\circ h:P_W\longrightarrow W
    \]
    is a projective cover of \(W\). In particular, \(\alpha\circ h\) is an essential epimorphism.
    Since \(h\) is surjective, \(\alpha\) is also surjective.
    
    Let \(U_0\leq U\) be a submodule such that
    \[
    \alpha(U_0)=W.
    \]
    Then
    \[
    (\alpha\circ h)\bigl(h^{-1}(U_0)\bigr)=W.
    \]
    Since \(\alpha\circ h\) is essential, we must have
    \[
    h^{-1}(U_0)=P_W.
    \]
    Applying \(h\), and using the surjectivity of \(h\), we get
    \[
    U=h(P_W)\subseteq U_0.
    \]
    Hence \(U_0=U\). Therefore \(\alpha:U\to W\) is an essential epimorphism.
    
    Finally, we show that in this situation
    \[
    h:P_W\longrightarrow U
    \]
    is a projective cover of \(U\). We already know that \(P_W\) is projective and that \(h\)
    is surjective. It remains to show that \(h\) is essential.
    
    Let \(X\leq P_W\) be a submodule such that
    \[
    h(X)=U.
    \]
    Then
    \[
    (\alpha\circ h)(X)=\alpha(U)=W.
    \]
    Since \(\alpha\circ h\) is essential, we obtain
    \[
    X=P_W.
    \]
    Thus \(h\) is an essential epimorphism. Therefore \(h:P_W\to U\) is a projective cover
    of \(U\).
    
    \medskip
    
    \textbf{(b)}
    Let
    \[
    q:U\longrightarrow U/V
    \]
    be the canonical epimorphism.
    
    First suppose that \(q\) is essential. We prove that
    \[
    V\subseteq \operatorname{Rad}U.
    \]
    If not, then there exists a maximal submodule \(M\) of \(U\) such that
    \[
    V\nsubseteq M.
    \]
    Since \(M\) is maximal and \(V\nsubseteq M\), we have
    \[
    M+V=U.
    \]
    Therefore
    \[
    q(M)=U/V.
    \]
    But \(M\neq U\), contradicting the essentiality of \(q\). Hence
    \[
    V\subseteq \operatorname{Rad}U.
    \]
    
    Conversely, suppose that
    \[
    V\subseteq \operatorname{Rad}U.
    \]
    Let \(X\leq U\) be a submodule such that
    \[
    q(X)=U/V.
    \]
    This is equivalent to
    \[
    X+V=U.
    \]
    We show that \(X=U\). Suppose, for contradiction, that \(X\neq U\). Since \(U\) is
    finite-dimensional, \(X\) is contained in some maximal submodule \(M\) of \(U\). Since
    \[
    V\subseteq \operatorname{Rad}U\subseteq M,
    \]
    we get
    \[
    U=X+V\subseteq M,
    \]
    which is impossible because \(M\neq U\). Therefore \(X=U\), and hence \(q\) is an
    essential epimorphism.
    
    Thus
    \[
    U\longrightarrow U/V
    \]
    is an essential epimorphism if and only if
    \[
    V\subseteq \operatorname{Rad}U.
    \]
    \end{proof}
    
    
    \begin{exercise}
    Let \(v\) be a valuation of a field \(K\). Prove the following properties:
    \begin{enumerate}
        \item[(i)] \(v(\pm 1)=0\).
        \item[(ii)] \(v(a)=v(-a)\).
        \item[(iii)] \(v(a^{-1})=-v(a)\).
        \item[(iv)] If \(v(b)<v(a)\), then
        \[
        v(a+b)=v(b).
        \]
    \end{enumerate}
    \end{exercise}
    
    \begin{proof}
    We use the defining properties of a valuation:
    \[
    v(xy)=v(x)+v(y),
    \qquad
    v(x+y)\geq \min\{v(x),v(y)\}.
    \]
    
    \textbf{(i)}
    Since \(1\cdot 1=1\), we have
    \[
    v(1)=v(1\cdot 1)=v(1)+v(1).
    \]
    Hence
    \[
    v(1)=0.
    \]
    Also,
    \[
    (-1)^2=1,
    \]
    so
    \[
    2v(-1)=v((-1)^2)=v(1)=0.
    \]
    Therefore
    \[
    v(-1)=0.
    \]
    Thus
    \[
    v(\pm 1)=0.
    \]
    
    \textbf{(ii)}
    For \(a\in K^\times\),
    \[
    -a=(-1)a.
    \]
    Therefore
    \[
    v(-a)=v((-1)a)=v(-1)+v(a)=v(a).
    \]
    
    \textbf{(iii)}
    For \(a\in K^\times\), we have
    \[
    aa^{-1}=1.
    \]
    Thus
    \[
    v(a)+v(a^{-1})=v(aa^{-1})=v(1)=0.
    \]
    Hence
    \[
    v(a^{-1})=-v(a).
    \]
    
    \textbf{(iv)}
    Assume that \(a,b\in K^\times\) and
    \[
    v(b)<v(a).
    \]
    First, \(a+b\neq 0\); otherwise \(a=-b\), and by part \textbf{(ii)} we would have
    \[
    v(a)=v(-b)=v(b),
    \]
    contradicting \(v(b)<v(a)\).
    
    By the valuation inequality,
    \[
    v(a+b)\geq \min\{v(a),v(b)\}=v(b).
    \]
    We now prove the reverse inequality. Since
    \[
    b=(a+b)-a,
    \]
    we have
    \[
    v(b)
    =
    v((a+b)-a)
    \geq
    \min\{v(a+b),v(-a)\}
    =
    \min\{v(a+b),v(a)\}.
    \]
    If \(v(a+b)>v(b)\), then both \(v(a+b)\) and \(v(a)\) are strictly larger than \(v(b)\),
    so
    \[
    \min\{v(a+b),v(a)\}>v(b),
    \]
    contradicting the previous inequality. Therefore
    \[
    v(a+b)\leq v(b).
    \]
    Combining this with
    \[
    v(a+b)\geq v(b),
    \]
    we obtain
    \[
    v(a+b)=v(b).
    \]
    \end{proof}

